\documentclass[aps,prb,onecolumn,nofootinbib,superscriptaddress]{revtex4-2}

\usepackage{amsmath,amssymb,amsthm,mathtools,bm}
\usepackage{booktabs}
\usepackage[colorlinks=true,linkcolor=blue,citecolor=blue,urlcolor=blue]{hyperref}
\usepackage{microtype}

\newcommand{\Tr}{\operatorname{Tr}}
\newcommand{\ii}{\mathrm{i}}
\newcommand{\cM}{\mathcal{M}}
\newcommand{\ket}[1]{\lvert #1\rangle}

\begin{document}

\title{\texorpdfstring{Local Gain--Loss Click Sequences at a Hard Controller Domain Wall}{Local Gain--Loss Click Sequences at a Hard Controller Domain Wall}}
\author{Asadullah Bhuiyan}
\affiliation{Department of Physics, Cornell University, Ithaca, New York 14853, USA}
\date{\today}

\begin{abstract}
We specify a trajectory-resolved pilot for local gain and loss records in a
number-conserving adaptive Gaussian circuit with a domain-wall-truncated
controller.  The simulation uses the canonical CPU dynamics entry point and
records, without covariance histories, the four measurement-feedback channels
generated by two Pauli trial eigenvectors and two flattened Hamiltonian bands.
We define a four-bit unit-cell motif, temporal and periodic spatial words,
shuffle null ensembles, and trajectory-bootstrap uncertainty.  Two matched
ten-trajectory schedules, a deterministic \(y\)-raster and a randomized site
order, separate schedule-conditioned patterns from candidates that persist
under update-order randomization.  This document is a methods specification:
it states estimators, validation gates, and interpretive limits but makes no
claim about the numerical outcome of the pilot.
\end{abstract}

\maketitle
\setcounter{tocdepth}{2}
\tableofcontents

\section{Gaussian instrument and trial eigenvectors}

The physical single-particle covariance is
\begin{equation}
 G=2C-\bm{1},\qquad
 C_{ij}=\Tr(\rho\,c_i^\dagger c_j),
\end{equation}
with two complex orbitals per unit cell.  At a controller center
\(R=(R_x,R_y)\), choose a Pauli operator
\(\sigma_{\mathrm{trial}}\).  Its normalized eigenvectors are denoted
\begin{equation}
 \sigma_{\mathrm{trial}}\tau_A=+\tau_A,\qquad
 \sigma_{\mathrm{trial}}\tau_B=-\tau_B.
 \label{eq:trial-eigenvectors}
\end{equation}
Thus \(A/B\) are the two eigenvectors of the chosen Pauli operator, not merely
post hoc channel names.  The pilot fixes
\(\sigma_{\mathrm{trial}}=\sigma_x\), or
\texttt{trial\_orbitals="X"}.  Projection of each trial eigenvector into the
upper and lower flattened-Hamiltonian bands gives four normalized observable
Wannier modes,
\begin{equation}
 \ket{w_{A+;R}},\quad \ket{w_{A-;R}},\quad
 \ket{w_{B+;R}},\quad \ket{w_{B-;R}}.
\end{equation}
The target occupations are
\begin{equation}
 s_{A+}=s_{B+}=0,\qquad s_{A-}=s_{B-}=1.
 \label{eq:targets}
\end{equation}
The implemented order inside a visited cell is \(A+,A-,B+,B-\).  Because a
measurement changes the covariance before the next channel is sampled, the
four outcomes at one cell are an ordered local instrument record.

\section{Hard-wall controller geometry}

The simulation fixes
\begin{equation}
 N_x=20,\quad N_y=24,\quad n_{\rm shell}=1,\quad
 \alpha_1=1,\quad \alpha_2=30,
 \label{eq:parameters}
\end{equation}
and constructs a topological slab using
\texttt{DW=True}.  The Wannier support is first cut to \(n_{\rm shell}=1\) and
then masked according to which side of the controller domain wall contains its
center:
\begin{equation}
 \ket{w^{\rm DW}_{a;R}}
 =
 \frac{M_R\ket{w_{a;R}}}
 {\lVert M_R\ket{w_{a;R}}\rVert}.
 \label{eq:wall-mask}
\end{equation}
The flags \texttt{dw\_truncation=True} and
\texttt{meas\_slab\_only=True} are mandatory.  For \(N_x=20\), the class rule
places the active slab at \(4\leq x\leq16\), with hard walls at \(x=4,16\).
The interface mask has width
\(n_{\rm shell}+1=2\), so the left and right interface columns are
\(\{4,5\}\) and \(\{15,16\}\); columns \(6,\ldots,14\) form the active
interior control.

Every trajectory is initialized with \texttt{init\_mode="maxmix"} and evolved
for 48 complete sweeps.  Cycles \(1,\ldots,24\) are burn-in and cycles
\(25,\ldots,48\) define the observation window.  No covariance history is
saved.  Production scripts call only
\begin{equation}
 \texttt{classA\_U1FGTN.run\_markov\_circuit(...)}
\end{equation}
with \texttt{G\_history=False} and \texttt{save=False}; the branch observer
streams the classical record before the final covariance is discarded.

\section{Signed click record and four-bit motifs}

For channel \(a\) at the \(m\)th visited cell, let
\(n_{a,m}\in\{0,1\}\) be the Born-rule measurement and
\(n'_{a,m}\in\{0,1\}\) the occupation after conditional feedback.  We retain
both the target defect
\begin{equation}
 X_{a,m}=\lvert n_{a,m}-s_a\rvert
\end{equation}
and signed particle transfer
\begin{equation}
 Y_{a,m}=n'_{a,m}-n_{a,m}\in\{-1,0,+1\}.
 \label{eq:signed-transfer}
\end{equation}
Negative, zero, and positive \(Y\) denote loss, no transfer, and gain.  Perfect
correction is used, hence
\begin{equation}
 X_{a,m}=\lvert Y_{a,m}\rvert
 \label{eq:perfect-identity}
\end{equation}
event by event.  Equation~\eqref{eq:targets} then restricts the four nonzero
click types as follows:
\begin{table}[ht]
\caption{Four click bits retained at each unit-cell visit.}
\begin{ruledtabular}
\begin{tabular}{cccc}
Bit & Controller channel & Trial eigenvector & Signed event \\
\midrule
0 & \(A+\) & \(\tau_A\), eigenvalue \(+1\) of \(\sigma_x\) & A-loss, \(Y=-1\) \\
1 & \(A-\) & \(\tau_A\), eigenvalue \(+1\) of \(\sigma_x\) & A-gain, \(Y=+1\) \\
2 & \(B+\) & \(\tau_B\), eigenvalue \(-1\) of \(\sigma_x\) & B-loss, \(Y=-1\) \\
3 & \(B-\) & \(\tau_B\), eigenvalue \(-1\) of \(\sigma_x\) & B-gain, \(Y=+1\) \\
\end{tabular}
\end{ruledtabular}
\end{table}

Define \(b_j\in\{0,1\}\) from these four events and encode the cell motif as
\begin{equation}
 q=b_0+2b_1+4b_2+8b_3\in\{0,\ldots,15\}.
 \label{eq:motif-code}
\end{equation}
This representation is deliberately not a single trinary A/B value.  A cell
can exhibit both A-loss and A-gain in one ordered visit; Eq.~\eqref{eq:motif-code}
preserves both bits instead of cancelling them.  Motif \(q=0\) is the
all-no-click state.

Each channel row stores the schedule, sample, one-based cycle, within-cycle
cell visit, zero-based global cell and channel steps, \(x,y\), interface label,
distances to both walls, \(A/B\) trial eigenvector, Pauli eigenvalue, band,
target, \(X\), \(Y\), click label, Born success probability, and validity.
A second table stores one row per unit-cell visit with \(q\) and all four bits.
The compressed array record remains the source of truth.

\section{Temporal and spatial word statistics}

For a fixed cell \(R\), temporal words of length \(k=2,3\) are
\begin{equation}
 W^{(t)}_{k}(R,t)=
 \big(q(R,t),q(R,t+1),\ldots,q(R,t+k-1)\big).
\end{equation}
Words never cross a trajectory, cell, or burn-in boundary.  At fixed
\((x,t)\), spatial words follow the positive periodic \(y\) direction:
\begin{equation}
 W^{(y)}_{k}(x,y,t)=
 \big(q(x,y,t),q(x,y+1,t),\ldots,q(x,y+k-1,t)\big),
 \qquad y\equiv y+N_y.
\end{equation}
Counts are reported for the left wall, right wall, pooled interface, active
interior, and entire active slab.

The temporal null independently permutes the post-burn-in cycle labels for
every trajectory and cell.  It therefore preserves the one-point motif count
of each cell while destroying temporal order.  The spatial null independently
permutes motif positions around each \((\text{trajectory},t,x)\) periodic
\(y\)-ring.  It preserves the ring's motif multiset while destroying
nearest-neighbor order.  Both null ensembles contain 1000 permutations
generated from a recorded analysis seed.

For observed word count \(N_w\), total count \(N\), and a word space of size
\(K=16^k\), the Jeffreys-regularized frequency is
\begin{equation}
 \widehat f_w=\frac{N_w+\tfrac12}{N+\tfrac12K}.
\end{equation}
The reported effect size is
\begin{equation}
 E_w=\log_2\frac{\widehat f_w}
 {\langle\widehat f_w^{\,\rm null}\rangle}.
 \label{eq:enrichment}
\end{equation}
Uncertainty is obtained from 1000 trajectory-level bootstrap resamples paired
with null-permutation replicates.  A ranked table retains words with at least
20 observations.  Its exploratory candidate flag only indicates that the
percentile interval for Eq.~\eqref{eq:enrichment} excludes zero; it is not a
multiple-testing-corrected discovery statement.

\section{Matched schedule design and deployment}

Two ten-trajectory ensembles use a common root seed \(20260814\) and identical
derived sample seeds.  The first uses \texttt{raster\_y}, which visits columns
left-to-right and sites within each column in increasing \(y\).  Only after
that campaign succeeds does the launcher begin the \texttt{random} schedule,
which reshuffles the active sites every cycle.  The schedules therefore form a
matched update-order comparison, not concurrently competing jobs.

Within a schedule, ten independent trajectories run on ten single-thread
process workers.  BLAS, OpenMP, MKL, and NumExpr thread counts are fixed to
one.  The launcher creates a new timestamped directory, refuses overwrites,
writes separate stage logs, and uses \texttt{\_SUCCESS} or
\texttt{\_FAILED} markers.  If the pre-existing CFT campaign remains active on
CPUs \(0\)--\(55\), the pilot is pinned to CPUs \(56\)--\(111\).  The
all-core topological-frustration campaign is allowed to finish before
production begins and is never interrupted.

\section{Validation gates}

A production stage is valid only if:
\begin{enumerate}
 \item every active cell is visited once per cycle and each visit contains
       exactly the four expected channel measurements;
 \item all valid Born success probabilities are finite and lie in \([0,1]\);
 \item \(A+,B+\) transfers belong to \(\{0,-1\}\), while \(A-,B-\) transfers
       belong to \(\{0,+1\}\);
 \item Eq.~\eqref{eq:perfect-identity} holds elementwise;
 \item dense-array dimensions agree with both tidy-table row counts;
 \item all geometry and controller parameters in Eq.~\eqref{eq:parameters},
       the Pauli basis, canonical dynamics entry point, seeds, and CPU policy
       are present in metadata;
 \item raster and random stages contain identical initialization and sample
       seed lists.
\end{enumerate}
Unit tests exhaust all 16 motif codes, simultaneous gain/loss events, temporal
boundary handling, periodic-\(y\) closure, shuffle count preservation,
determinism, and malformed records.  A \(4\times6\), two-sample, two-cycle run
validates both schedules and the full artifact schema before production.

\section{Interpretive limits}

This is a ten-trajectory pilot designed to identify candidate effects, not
thermodynamic or long-time laws.  A raster-only spatial word may be generated
by the update schedule itself.  Persistence under random ordering is a
stronger robustness criterion, but does not by itself establish universality.
The hard mask is a boundary of the controller support, not a smooth material
interface.  Perfect correction also identifies wrong outcomes with absolute
transfer activity and therefore does not test failed feedback.  These limits
must accompany every numerical summary derived from the pilot.

\end{document}
